\documentclass[10pt,a4paper]{article}

\usepackage{microtype}
\usepackage{a4wide}
\usepackage{tgpagella}
\usepackage[euler-digits]{eulervm}
\usepackage{mdwlist}
\usepackage[T1]{fontenc}
\usepackage{isabelle,isabellesym}

\usepackage{amssymb}
\usepackage{amsmath}
\usepackage{eurosym}
\usepackage[only,bigsqcap,bigparallel,fatsemi,interleave,sslash]{stmaryrd}
\usepackage{eufrak}
\usepackage{textcomp}
\usepackage{wasysym}

\usepackage{tikz}
\usetikzlibrary{arrows.meta,positioning,calc}

\usepackage{xcolor}
\definecolor{mlkblue}{HTML}{2563EB}
\definecolor{mlklightblue}{HTML}{DBEAFE}

% this should be the last package used
\usepackage{pdfsetup}

% urls in roman style, theory text in math-similar italics
\urlstyle{rm}
\isabellestyle{it}

% for uniform font size
\renewcommand{\isastyle}{\isastyleminor}

% suppress proof bodies, keep theorem statements
\isadroptag{proof}

\renewcommand{\isasymlonglonglongrightarrow}{\longlonglongrightarrow}

\begin{document}

\title{ML-KEM Functional Correctness in Isabelle/HOL}
\author{mlkem-native contributors}
\maketitle

\begin{abstract}
Machine-checked functional correctness proofs for the polynomial
arithmetic routines in \texttt{mlkem-native}, an ML-KEM (FIPS~203)
implementation in portable C\@.  The C code is automatically translated
to Isabelle/HOL via \emph{AutoCorrode}, and each function is verified
against an abstract integer-level specification.  The Number Theoretic
Transform (NTT) and its inverse are further shown to be two-sided
inverses modulo~$q$.
\end{abstract}

\tableofcontents

% sane default for proof documents
\parindent 0pt\parskip 0.5ex

% generated text of all theories
\input{session}

\end{document}

%%% Local Variables:
%%% mode: latex
%%% TeX-master: t
%%% End:
